% ============================================
% 第1章：引言 + 任务概述 + 数据描述
% 负责人：A
% ============================================

\section{引言}

\subsection{编写目的}

本文档为在线论文文档自动生成工具（Thesis Generator）软件需求规格说明书，用于完整定义系统的全部功能性与非功能性需求，明确项目开发边界、业务规则、数据规范和运行约束。

本文档的预期读者包括：
\begin{itemize}
    \item 项目经理：用于制定项目计划和跟踪需求实现情况；
    \item 系统设计人员：作为概要设计和详细设计的技术依据；
    \item 前后端开发工程师：明确各自负责模块的功能边界和接口约定；
    \item 测试人员：作为编写测试用例和验收测试的标准参考。
\end{itemize}

\subsection{项目背景}

本项目为高校大二软件工程实训项目（CST328 项目实训），开发单位为在校实训开发小组，共六名成员。项目目标是开发一套Web端在线论文编辑与标准化文档生成系统，服务于高校学生在课程论文、毕业论文、项目设计类论文撰写过程中的规范化需求。

当前高校学生在撰写论文时面临以下痛点：论文模板格式不统一、排版操作繁琐、参考文献格式手动调整耗时、文档导出格式错乱、课程论文批阅流程缺乏线上化支持。本系统旨在通过Web平台实现从模板分发、在线写作、预览核查、文档导出到教师线上批改的一体化论文生成工具，降低学生论文撰写的格式负担，提高教师批阅效率。

系统独立部署运行，无第三方外部系统强制对接，仅依托Web服务、关系型数据库、Redis缓存和文件存储完成完整业务闭环。

\subsection{术语与定义}

\begin{itemize}
    \item \textbf{论文模板}：预设的论文结构框架，包含封面字段配置、章节结构、排版格式参数（字体、字号、行距、页边距等），学生创建论文时自动套用；
    \item \textbf{模板版本}：同一模板的多次修订记录，每次修改生成新版本，学生新建论文自动使用最新启用版本；
    \item \textbf{富文本编辑器}：基于Web浏览器的所见即所得（WYSIWYG）在线编辑器，支持文字格式化、图片嵌入、表格插入、引用标记等功能；
    \item \textbf{批注}：教师在学生论文指定文字位置添加的修改意见或评论，与论文原文位置绑定；
    \item \textbf{批阅}：教师对学生提交的课程论文进行审查、添加批注、评定分数和等级的操作流程；
    \item \textbf{提交}：学生完成论文编辑后将论文发送至教师端等待批阅的操作，提交后论文锁定不可编辑；
    \item \textbf{退回}：教师将不符合要求的论文返回给学生重新修改的操作，退回后论文解锁恢复编辑；
    \item \textbf{DOCX/PDF导出}：将学生在线编辑的论文内容按照模板格式配置生成标准Word文档或PDF文件的功能。
\end{itemize}

\subsection{参考资料}

\begin{enumerate}
    \item GB/T 9385-2008 计算机软件需求规格说明规范；
    \item 2026软件工程实训项目任务书；
    \item Web前后端分离开发设计规范；
    \item Apache POI 官方文档——Java操作Microsoft Office文档的开源库；
    \item GB/T 7714-2015 信息与文献 参考文献著录规则；
    \item Spring Boot 官方参考文档；
    \item PostgreSQL 数据库设计开发规范。
\end{enumerate}

\section{任务概述}

\subsection{目标}

搭建一套Web端在线论文编辑与标准化文档生成平台，区分系统管理员、学生、教师三类用户角色，分别赋予模板管理、论文编辑导出、课程论文线上批阅等对应权限，实现以下核心业务目标：

\begin{enumerate}
    \item 管理员端：支持按学院维度对论文模板进行分类管理，可自定义发布多类型模板（毕业论文、课程论文、项目论文），模板内置封面、中英文摘要、多级章节正文、参考文献等固定结构，统一约束字体、行距、页边距等排版格式，支持模板版本管理与在线预览；
    \item 学生端：可按需选用模板，依托结构化章节大纲在线完成论文内容编写（含文字排版、图片嵌入、表格编辑、参考文献管理），支持云端草稿存储、全文Web预览、一键导出标准DOCX和PDF格式文档，课程论文可提交至教师端等待批阅；
    \item 教师端：可查看学生提交的课程论文内容，在文档指定位置添加批注修改意见，填写整体评语并评定分数与等级，支持退回论文供学生修改重提，留存历次批阅历史记录。
\end{enumerate}

\subsection{运行环境}

\begin{itemize}
    \item \textbf{服务端操作系统}：Windows Server 2019 / Linux CentOS 7 及以上；
    \item \textbf{Web服务容器}：Spring Boot内嵌Tomcat / Nginx反向代理；
    \item \textbf{数据库}：PostgreSQL 16+，存储全部业务数据；
    \item \textbf{缓存与队列}：Redis 7+，用于Token缓存、热点数据缓存、异步任务队列；
    \item \textbf{文件存储}：服务器本地文件系统，存储上传的论文图片资源及生成的DOCX/PDF文件；
    \item \textbf{客户端环境}：Chrome（最新版）、Edge（最新版）、Firefox（最新版）等主流现代浏览器。
\end{itemize}

\subsection{条件与限制}

\begin{enumerate}
    \item 系统仅开放管理员、学生、教师三类角色，不开放游客的论文编辑和导出权限；
    \item 用户登录采用JWT Token鉴权机制，所有API接口需携带有效Token，按角色严格区分数据访问权限；
    \item 学生论文编辑过程中上传的图片单张不得超过5MB，格式限制为jpg/png/gif；
    \item DOCX和PDF文档生成过程涉及服务端文件IO操作，单次导出操作应在60秒内完成；
    \item 系统为独立部署的Web应用，不依赖第三方外部系统接口，所有功能在封闭环境中完成；
    \item 服务器硬件资源有限（实训环境），需合理控制并发导出任务数量和文件存储空间。
\end{enumerate}

\section{数据描述}

\subsection{静态数据}

系统长期持久存储且不频繁变更的基础数据包括：

\begin{itemize}
    \item \textbf{用户账号信息}：用户名、加密密码、真实姓名、角色标识（ADMIN/STUDENT/TEACHER）、所属学院；
    \item \textbf{学院基础信息}：学院名称、学院编号；
    \item \textbf{论文模板信息}：模板名称、论文类型（毕业论文/课程论文/项目论文）、关联学院、启用状态；
    \item \textbf{模板版本配置}：封面字段配置（JSON）、章节结构定义（JSON）、排版格式参数（JSON）；
    \item \textbf{题目讨论与公告}：系统公告内容、讨论帖子内容（本系统暂不涉及）。
\end{itemize}

\subsection{动态数据}

系统运行过程中实时产生、频繁变更的数据包括：

\begin{itemize}
    \item \textbf{用户登录会话}：JWT Token及其Redis缓存，具有时效性（24小时过期）；
    \item \textbf{论文内容数据}：学生在线编辑的论文章节HTML内容、草稿自动保存版本；
    \item \textbf{参考文献条目}：学生逐条录入的参考文献信息（作者、标题、刊物、年份等）；
    \item \textbf{上传图片记录}：学生在论文中嵌入的图片文件及元数据；
    \item \textbf{提交记录}：学生历次提交论文的时间戳和版本号；
    \item \textbf{批注与批阅记录}：教师在论文指定位置添加的批注意见、整体评语、分数和等级；
    \item \textbf{导出任务状态}：DOCX/PDF生成任务的进度和结果。
\end{itemize}

\subsection{数据库介绍}

采用PostgreSQL关系型数据库存储全部业务数据。数据库名称为\verb|thesis_generator|，设计包含以下核心数据表：

\begin{itemize}
    \item \verb|sys_user|：系统用户表，存储全部用户账号、角色和所属学院信息；
    \item \verb|sys_college|：学院信息表，存储学院基础数据；
    \item \verb|template|：论文模板表，存储模板基本信息和启用状态；
    \item \verb|template_version|：模板版本表，存储每个模板的各版本配置数据（封面字段、章节结构、格式参数），使用JSONB类型字段以支持灵活的配置结构；
    \item \verb|thesis|：论文主表，存储每篇论文的元信息（标题、状态、关联学生和模板版本）；
    \item \verb|thesis_section|：论文章节表，存储每篇论文各章节的结构化信息和HTML内容；
    \item \verb|thesis_reference|：参考文献表，存储每篇论文的参考文献条目；
    \item \verb|thesis_image|：图片资源表，记录学生在论文中上传的图片文件信息；
    \item \verb|submission|：提交记录表，记录学生历次提交论文的时间信息；
    \item \verb|annotation|：批注表，存储教师在论文指定位置添加的批注意见；
    \item \verb|review_record|：批阅记录表，存储教师的评语、分数、等级和操作类型。
\end{itemize}

各数据表之间通过外键建立关联约束：用户关联学院、论文关联用户和模板版本、章节和参考文献关联论文、批注和批阅记录关联论文等。

\subsection{数据词典}

以下是系统核心数据元素的简要词典：

\begin{table}[htbp]
    \centering
    \caption{核心数据词典}
    \begin{tabular}{|c|c|c|c|}
        \hline
        \textbf{数据项} & \textbf{所属表} & \textbf{数据类型} & \textbf{说明} \\
        \hline
        用户ID & sys\_user & BIGINT & 自增主键，唯一标识用户 \\
        \hline
        用户名 & sys\_user & VARCHAR(50) & 登录账号，唯一 \\
        \hline
        角色 & sys\_user & VARCHAR(20) & ADMIN/STUDENT/TEACHER \\
        \hline
        论文状态 & thesis & VARCHAR(20) & DRAFT/COMPLETED/SUBMITTED/REVIEWED/RETURNED \\
        \hline
        章节内容 & thesis\_section & TEXT & HTML格式的章节正文 \\
        \hline
        模板类型 & template & VARCHAR(20) & GRADUATION/COURSE/PROJECT \\
        \hline
        封面字段配置 & template\_version & JSONB & 定义封面显示哪些字段及其排序 \\
        \hline
        批注偏移量 & annotation & INT & 选中文字在章节HTML中的起始字符位置 \\
        \hline
        分数 & review\_record & INT & 0-100整数，教师评定 \\
        \hline
        等级 & review\_record & VARCHAR(10) & 自动换算：$\geq$90优/≥80良/≥70中/≥60及格 \\
        \hline
    \end{tabular}
\end{table}

\subsection{数据采集}

\begin{itemize}
    \item \textbf{用户数据}：学生自主注册（选择学院、设定账号密码）、教师由管理员后台录入或自主注册；
    \item \textbf{模板数据}：管理员在后台管理界面手动配置封面字段、章节结构和格式参数；
    \item \textbf{论文内容数据}：学生在Web编辑器中实时编辑，通过自动保存（每30秒）和手动保存两种方式持久化到数据库；
    \item \textbf{参考文献数据}：学生通过专用表单逐条录入，支持BibTeX格式批量导入；
    \item \textbf{图片数据}：学生在前端通过文件选择器上传，服务端校验格式和大小后存储；
    \item \textbf{批阅数据}：教师在批阅界面选中论文文字添加批注，填写评语和分数后提交；
    \item \textbf{导出文件}：学生点击导出按钮触发服务端异步生成任务，完成后提供下载。
\end{itemize}
